\subsection{EXPONENTIAL IN COMPLETE NORMED ASSOCIATIVE ALGEBRAS}

In this no. we denote by A a complete normed unital associative algebra (General Topology, Chapter IX, § 3, no. 7). Then \(\| x.y\| \leqslant \| x\|\) . \(\| y\|\) for \(x,y\) in A.

Let I be a finite set and let \(\hat{\mathrm{P}} (\mathrm{A}^{\mathrm{I}};\mathrm{A})\) be the vector space of formal power series with continuous components on \(\mathbf{A}^{\mathrm{I}}\) with values in A (Differentiable and Analytic Manifolds, R, Appendix, no. 5) with the algebra structure obtained by writing

\[
f. g = m \circ (f, g) \quad \text { for } f, g \text { in } \hat {\mathrm{P}} (\mathrm{A} ^ {\mathrm{I}}; \mathrm{A}),
\]

where \(m: \mathbf{A} \times \mathbf{A} \to \mathbf{A}\) denotes multiplication on A. Arguing as in no. 1 and using Proposition 1 of § 5, no. 1, we define a continuous homomorphism of unital algebras \(u \mapsto \tilde{u}\) of \(\hat{\mathrm{A}}(\mathrm{I})\) into \(\hat{\mathrm{P}}(\mathrm{A}^{\mathrm{I}}; \mathrm{A})\) mapping the indeterminate of index \(i\) to \(\mathbf{pr}_{\mathrm{i}}\) ; this homomorphism extends the Lie algebra homomorphism of \(\hat{\mathrm{L}}(\mathrm{I})\) into \(\hat{\mathrm{P}}(\mathrm{A}^{\mathrm{I}}; \mathrm{A})\) defined in no. 1. If \(u = \sum_{\nu} u_{\nu}\) with \(u_{\nu} \in \mathrm{A}^{\nu}(\mathrm{I})\) for \(\nu \in \mathbf{N}^{\mathrm{I}}\) , then \(\tilde{u} = \sum_{\nu} \tilde{u}_{\nu}\) , where \(\tilde{u}_{\nu}\) is the polynomial mapping \((t_{i})_{i \in I} \mapsto u_{\nu}((t_{i}))\) .

Let \(\boldsymbol{u} = (u_{j})_{j \in I}\) be a finite family of elements of \(\hat{\mathbf{A}}(\mathbf{I})\) , let \(v \in \hat{\mathbf{A}}(\mathbf{J})\) and write \(w = v \circ \boldsymbol{u}\) ( \(\S 5\) , no. 1). Then

\[
(v \circ \boldsymbol {u}) ^ {\sim} = \tilde {v} \circ \tilde {\boldsymbol {u}}.\tag{18}
\]

This follows by extending by continuity formula (2) of § 5, no. 1 and from Differentiable and Analytic Manifolds, R, Appendix, no. 6.

In particular we take \(\mathrm{I} = \{\mathrm{U}\}\) , identify A and \(\mathrm{A}^{(\mathrm{U})}\) and consider the images \(\tilde{e}\) and \(\tilde{l}\) of the series \(e(\mathrm{U}) = \sum_{n\geqslant 1}\mathrm{U}^n /n!\) and \(l(\mathrm{U}) = \sum_{n\geqslant 1}(-1)^{n - 1}\mathrm{U}^n /n\) in \(\hat{\mathrm{P}} (\mathrm{A};\mathrm{A})\) . Then \(\| \tilde{\mathbf{U}}^n\| \leqslant 1\) for \(\| x_1\dots x_n\| \leqslant \| x_1\| \dots \| x_n\|\) for \(x_{1},\ldots ,x_{n}\) in A. Therefore the radius of absolute convergence of \(\tilde{e}\) (resp. \(\tilde{l}\) ) is infinite (resp. \(\geqslant 1\) ).

We shall denote by \(e_{A}\) (resp. \(l_{A}\) ) the analytic mapping of A into A (resp. of B into A, where B is the open unit ball of A) defined by the convergent series \(\tilde{e}\) (resp. \(\tilde{l}\) ) and we shall write \(\exp_{\mathbf{A}}(x)=1+e_{\mathbf{A}}(x)\) (for \(x\in A\) ) and

\[
\log_ {\mathrm{A}} (x) = l _ {\mathrm{A}} (x - 1)
\]

(for \(x \in \mathbf{A}\) , \(\| x - 1\| < 1\) ). Then

\[
\exp_ {\mathrm{A}} x = \sum_ {n \geqslant 0} \frac {x ^ {n}}{n !} (x \in \mathrm{A})\tag{19}
\]

\[
\log_ {\mathrm{A}} x = \sum_ {n \geqslant 1} (- 1) ^ {n - 1} \frac {(x - 1) ^ {n}}{n} (x \in \mathrm{A}, \| x - 1 \| <   1).\tag{20}
\]

As \((e\circ l)(\mathrm{U}) = (l\circ e)(\mathrm{U}) = \mathrm{U}\) (cf.~§ 6, no. 1), by (18) \(\tilde{e}\circ \tilde{l} = \tilde{l}\circ \tilde{e} = \mathrm{Id}_{\mathbf{A}}\) Therefore (Differentiable and Analytic Manifolds, R, 3.1.9)

\begin{enumerate}
\def\labelenumi{(\arabic{enumi})}
\setcounter{enumi}{20}
\tightlist
\item
\end{enumerate}

\[
\exp_ {\mathrm{A}} (\log_ {\mathrm{A}} (x)) = x \quad (x \in \mathrm{A}, \| x - 1 \| \leqslant 1)\tag{21}
\]

\[
\log_ {\mathbf {A}} (\exp_ {\mathbf {A}} (x)) = x \quad (x \in \mathrm{A}, \| x \| <   \log 2)\tag{22}
\]

for \(\| x\| < \log 2\) implies \(\| \exp_{\mathbf{A}}(x) - 1\| \leqslant \exp \| x\| -1 < 1.\)

Finally we consider A as a complete normed Lie algebra. Then

\[
\| [ x, y ] \| = \| x y - y x \| \leqslant 2 \| x \|. \| y \|.
\]

Proposition 1 of no. 2 implies that the domain of absolute convergence of the formal power series \(\tilde{\mathbf{H}}\) contains the set

\[
\Omega = \{(x, y) \in \mathrm{A} \times \mathrm{A} | \| x \| + \| y \| <   \frac {1}{2} \log 2 \}.
\]

Hence \(\tilde{\mathbf{H}}\) defines an analytic function \(h\colon \Omega \to \mathrm{A}\) . Then \(h(x,y) = \sum_{r,s}\mathrm{H}_{r,s}(x,y)\) (cf.~§ 3, no. 1, Remark 4).

PROPOSITION 3. For \(\| x\| + \| y\| < \frac{1}{2}\log 2,\)

\[
\exp_ {A} x. \exp_ {A} y = \exp_ {A} h (x, y).\tag{23}
\]

It follows from (18) and the relation \(e^{\mathrm{U}}e^{\mathrm{V}} = e^{\mathrm{H}(\mathrm{U},\mathrm{V})}\) that

\[
m \circ (1 + \tilde {e}, 1 + \tilde {e}) = (1 + \tilde {e}) \circ \tilde {\mathrm{H}}
\]

in \(\hat{\mathrm{P}} (\mathrm{A}\times \mathrm{A};\mathrm{A})\) . We therefore deduce from Differentiable and Analytic Manifolds, R, 3.1.9 that (23) is true for \((x,y)\) sufficiently close to \((0,0)\) , whence the proposition follows by analytic continuation (Differentiable and Analytic Manifolds, R, 3.2.5).

